\section{Исторический контекст и значение дипломных работ}
\label{sec:history}

Выпускная квалификационная работа (ВКР) является традиционным итогом высшего образования, однако в XXI веке её роль претерпела значительные изменения. Если ранее диплом часто воспринимался как формальность, то сегодня он превратился в полноценный научно-исследовательский или прикладной проект, способный внести реальный вклад в индустрию.

\subsection{Эволюция требований к ВКР}
\label{subsec:evolution}

Согласно исследованию \cite{edutrends}, за последние два десятилетия объём и сложность дипломных работ возросли в среднем на 40\%. Это связано с развитием технологий и повышением требований работодателей к выпускникам. В таблице \ref{tab:comparison} приведены сравнительные характеристики дипломных проектов разных лет.

\begin{table}[htbp]
	\centering
	\caption{Сравнение подходов к выполнению дипломных работ}
	\label{tab:comparison}
	\begin{tabular}{|p{4cm}|p{4.5cm}|p{4.5cm}|}
		\hline
		\textbf{Параметр} & \textbf{1990–2000 гг.} & \textbf{2010–2025 гг.} \\
		\hline
		Преобладающий тип & Теоретический обзор & Разработка прототипа / исследование \\
		\hline
		Объём кода (для IT-направлений) & До 1000 строк & 5000–20000 строк \\
		\hline
		Использование систем контроля версий & Отсутствует & Обязательно (Git) \\
		\hline
		Внешнее рецензирование & Факультативно & Обязательно \\
		\hline
	\end{tabular}
\end{table}

\subsection{Пример оформления листинга}
Для демонстрации возможностей шаблона приведём фрагмент кода на языке C с использованием специальных маркеров для git-нотации (подсветка добавленных и удалённых строк).

\begin{lstlisting}[caption={Пример кода с git-стилем}, label=lst:example]
	#include <stdio.h>
	
	int main() {
		printf("Hello, World!\n");
@--      // old version -@
@--      printf("Old message\n"); -@
@--       -@
@++    // new version +@
@++    printf("Updated message\n"); +@
@++   +@
		return 0;
	}
\end{lstlisting}

Как видно, шаблон позволяет выделять изменения цветом: удалённые строки окрашены в красный, добавленные — в зелёный.

\subsection{Роль диплома в академической карьере}
Многие исследователи \cite{ispras, career} отмечают, что качественная дипломная работа часто становится первым серьёзным достижением молодого специалиста, открывая путь в магистратуру, аспирантуру или в ведущие IT-компании. На рисунке \ref{fig:career} схематично изображены возможные траектории после защиты.

\begin{figure}[htbp]
	\centering
	\includegraphics[width=0.7\textwidth]{example-image} % для демонстрации можно использовать стандартный график
	\caption{Возможные карьерные траектории выпускника}
	\label{fig:career}
\end{figure}
(В реальном документе здесь должна быть размещена диаграмма, построенная средствами \texttt{tikz} или внешним изображением.)

Таким образом, исторически сложилось, что дипломная работа остаётся ключевым этапом оценки компетенций выпускника, а современные тенденции лишь усиливают её значимость.